\subsection{Le problème du tri}

\textbf{Entrée :} une séquence de $n$ nombres $\langle a_1, a_2, \dots, a_n \rangle$.\\
\textbf{Sortie :} une permutation $\langle a'_1, a'_2, \dots, a'_n \rangle$ telle que
$a'_1 \le a'_2 \le \dots \le a'_n$.

% ─────────────────────────────────────────────────────────────────────────────
\subsection{Tri par insertion (\textit{Insertion Sort})}

L'idée est celle du tri d'une main de cartes : on parcourt les éléments un par un et on
insère chaque nouvel élément à sa place correcte parmi ceux déjà triés.

\begin{center}
\begin{tikzpicture}[font=\small]
  % Array cells
  \foreach \val/\idx in {5/1,2/2,4/3,6/4,1/5,3/6}{
    \node[draw, minimum width=0.9cm, minimum height=0.7cm,
          fill=blue!8] at (\idx*1.0-0.5, 0) {\val};
    \node[font=\tiny, gray] at (\idx*1.0-0.5, -0.55) {\idx};
  }
  % Arrow showing "key"
  \draw[->, keyred, thick] (2.5, -0.9) -- (2.5, -0.4)
        node[midway, right=2pt, keyred] {\textit{key}};
  % Already sorted bracket
  \draw[decorate, decoration={brace,amplitude=5pt}, keygreen]
        (0.05, 0.38) -- (1.45, 0.38)
        node[midway, above=5pt, keygreen] {trié};
\end{tikzpicture}
\end{center}

\begin{algorithm}[H]
\caption{Insertion-Sort($A$)}
\begin{algorithmic}[1]
\For{$j = 2$ \textbf{to} $A.length$}
    \State $key \leftarrow A[j]$  \AlgComment{élément à insérer}
    \State $i \leftarrow j - 1$
    \While{$i > 0$ \textbf{and} $A[i] > key$}
        \State $A[i+1] \leftarrow A[i]$  \AlgComment{décaler vers la droite}
        \State $i \leftarrow i - 1$
    \EndWhile
    \State $A[i+1] \leftarrow key$  \AlgComment{insérer à la bonne place}
\EndFor
\end{algorithmic}
\end{algorithm}

% ─────────────────────────────────────────────────────────────────────────────
\subsection{Preuve de correction : invariant de boucle}

Pour prouver qu'un algorithme itératif est correct, on utilise un \textbf{invariant de
boucle}, une propriété vérifiée à trois moments clés :

\begin{enumerate}
  \item \textbf{Initialisation} : vraie avant la première itération.
  \item \textbf{Conservation} : si vraie avant une itération, reste vraie avant la suivante.
  \item \textbf{Terminaison} : à la fin de la boucle, l'invariant prouve la correction.
\end{enumerate}

\textbf{Invariant pour le tri par insertion :} au début de chaque itération $j$, le
sous-tableau $A[1 \ldots j-1]$ contient les éléments initialement en $A[1 \ldots j-1]$,
\emph{mais triés}.

% ─────────────────────────────────────────────────────────────────────────────
\subsection{Analyse de complexité (modèle RAM)}

On utilise le modèle \textbf{Random-Access Machine (RAM)} : chaque instruction simple
(arithmétique, accès mémoire, test) prend un temps constant.

\begin{center}
\renewcommand{\arraystretch}{1.3}
\begin{tabular}{lll}
\toprule
\textbf{Cas} & \textbf{Condition} & \textbf{Complexité} \\
\midrule
Meilleur cas & Tableau déjà trié (while ne s'exécute jamais) & $\Theta(n)$ \\
Pire cas     & Tableau trié en ordre inverse & $\Theta(n^2)$ \\
\bottomrule
\end{tabular}
\end{center}

\textbf{Calcul du pire cas :}
\[
  T(n) = \sum_{j=2}^{n} (j-1) = \frac{n(n-1)}{2} = \Theta(n^2)
\]

\begin{tcolorbox}[colback=lightblue, colframe=keyblue, fonttitle=\bfseries,
                  title={Pourquoi se concentrer sur le pire cas ?}]
Il fournit une \textbf{garantie supérieure} sur le temps d'exécution, valable pour toute
entrée possible. C'est la référence standard en analyse d'algorithmes.
\end{tcolorbox}